% Part: first-order-logic
% Chapter: introduction
% Section: syntax

\documentclass[../../../include/open-logic-section]{subfiles}

\begin{document}

\olfileid{fol}{int}{syn}

\olsection{Sintaksis}

Sepisan, kita kudu njlentrehake kanthi trep rerangken simbol endi kang
dianggep !!{sentence}s ing logika sepisanan. Bab iki bakal ditindakake
mengko; sauntara iki kita bakal maju nganggo tuladha. Perangan pambangun
dhasar---tetembungane---ing logika sepisanan kabagi dadi rong perangan.
Perangan kapisan yaiku simbol-simbol kang digunakake kanggo ngandharake
bab tartamtu utawa nuduhake sawijining bab kanthi mligi. Kita nuduhake
sawijining bab nganggo !!{constant}s, lan ngandharake bab ngenani apa
kang dituduhake mau nganggo !!{predicate}s. Tuladhane, kita bisa
nggunakake $\Obj a$ minangka !!a{constant} kanggo nuduhake siji bab,
banjur ngandharake sawijining bab ngenani bab kasebut nganggo
!!{sentence}~$\Atom{\Obj P}{\Obj a}$. Yen sampeyan wis duwe makna kanggo
``$\Obj a$'' lan ``$\Obj P$'' ing pikiran, sampeyan bisa maca
$\Atom{\Obj P}{\Obj a}$ minangka ukara basa Inggris (lan mbokmenawa
pancen wis nindakake mangkono nalika sepisanan sinau logika formal).
Sawise nduweni !!{sentence}s logika sepisanan kang prasaja kaya
mangkono, kita bisa mbangun ukara kang luwih majemuk nganggo perangan
kapindho saka tetembungan kasebut, yaiku simbol logis (konektif lan
kuantor). Dadi, tuladhane, kita bisa mbentuk ekspresi kaya
$(\Atom{\Obj P}{\Obj a} \land \Atom{\Obj Q}{\Obj b})$
utawa~$\lexists[\Obj x][\Atom{\Obj P}{\Obj x}]$.

Supaya bisa menehi definisi semantik lan aturan sistem !!{derivation}
kang trep lan dibutuhake kanggo panliten metalogis kang tliti, sepisan
kita kudu menehi definisi kang trep ngenani apa kang dianggep
!!a{sentence} ing logika sepisanan. Gagasan dhasare cukup gampang
dimangerteni: ana sawetara !!{sentence}s prasaja kang bisa kita bentuk
mung saka !!{predicate}s lan !!{constant}s, kayata~$\Atom{\Obj P}{\Obj
a}$. Sabanjure, saka ukara-ukara kasebut kita mbentuk ukara-ukara kang
luwih majemuk nganggo konektif lan kuantor. Nanging, aturan kang
pas endi sing marengake kita mbentuk !!{sentence}s kang luwih majemuk?
Aturan kasebut kudu ditemtokake; yen ora, kita durung nemtokake
``!!{sentence} ing logika sepisanan'' kanthi cukup trep. Ana sawetara
prekara. Prekara kapisan yaiku nemtokake rerangken simbol kang trep
supaya dianggep !!{sentence}s. Prekara kapindho yaiku nindakake iki
kanthi cara kang ndadekake kita bisa menehi bukti matematis ngenani
\emph{kabeh} !!{sentence}s. Pungkasane, kita uga kudu menehi definisi
kang trep kanggo sawetara operasi dhasar tumrap !!{sentence}s, kayata
``ganti saben $\Obj x$ ing~$!A$ nganggo~$\Obj b$.'' Masalahe, kuantor
lan !!{variable}s kang ana ing logika sepisanan ndadekake cara
nindakake iki ora mesthi cetha. Tuladhane, apa $\lexists[\Obj
x][\Atom{\Obj P}{\Obj a}]$ kudu dianggep !!a{sentence}? Kepriye karo
$\lexists[\Obj x][\lexists[\Obj x][\Atom{\Obj P}{\Obj x}]]$? Apa asil
saka ``ganti $\Obj x$ nganggo~$\Obj b$ ing $(\Atom{\Obj
P}{\Obj x} \land \lexists[\Obj x][\Atom{\Obj P}{\Obj x}])$''?

\end{document}
